\begin{textsample}{POS Dim 2 – vem\_ed\_23 – Score 1.00 – vem\_ed\_23/t123.txt}  \label{ex:f2_pos_032}
Quem é Quem Cathy Lee T.
Arcuino é diretora executiva de Global Engagement (em tradução livre, engajamento para a cidadania global) na Wilkes University (EUA).
Logo após sua graduação, em 1999, atuou como voluntária do Peace Corps em uma pequena vila no sul do Cazaquistão, próximo à fronteira uzbeque.
Ensinou inglês no Japão, na Tailândia e na Polônia, trabalhou numa ONG no Quirguistão e participou do programa Fulbright International Educator Administrator no Japão.
Mestre e doutora em Educação (respectivamente, pela Framingham State College e pela Colorado State University), é formada pela NAFSA Academy (Associação de Educadores Internacionais).
Conte como suas origens influenciaram sua visão de cidadania global.
Meus pais são filipinos que emigraram para os Estados Unidos.
Nasci em Nova York e quando tinha 4 anos nos mudamos para a Califórnia.
Meus pais me ensinaram que eu teria que trabalhar o dobro, não só porque sou uma mulher nos Estados Unidos, mas por ser asiática.
Tenho doutorado em Educação e meus pais ficaram felizes porque queriam que eu fosse “doutora” médica, mas descobriram que me tornara Dra.
Arcuino.
Na juventude, visitava minha família nas Filipinas e observava os recursos educacionais que tive e meus primos não tiveram lá.
E me perguntava: e se meus pais não tivessem emigrado?
Como percebo que esse pequeno movimento que meus pais fizeram mudou completamente minha vida, quero transformar outras vidas por meio da internacionalização.
Fale sobre sua experiência no Peace Corps.
Quando ingressei como voluntária no Peace Corps, minha expectativa era ir para as Filipinas.
Só que me disseram que teria de aguardar um ano e não queria esperar tanto, então fui parar no Cazaquistão.
Para ser honesta, tinha 21 anos e não fazia ideia de onde ficava o país.
Foi uma das melhores decisões que fiz porque conheci um país e uma cultura lindas.
Também aprendi muito sobre mim mesma.
Não falava nada de russo, só sabia que “niet” significa “não”.
Naquela época o russo era a língua principal, devido à influência soviética, especialmente no sul do país, onde vivi, próximo à fronteira uzbeque.
Como foi sua experiência como professora de inglês no Japão?
Foi um grande um choque cultural.
Logo depois de viver numa pequena vila cazaque, fui para Osaka, Japão.
Também foi muito diferente porque naquela época [ 2001 ] o Japão estava muito avançado no ensino-aprendizagem online.
Eu vivia em Osaka, mas meus alunos eram de todo o país.
Era exatamente como estamos nos encontrando agora.
Não tinha mais do que três estudantes online e eu ensinava inglês.
Nunca poderia imaginar que vinte e poucos anos depois isso seria parte do mundo, porque naquela época isso já estava acontecendo no Japão.
No Japão se trabalha muitas \textbf{horas} por dia, então para fazer cursos extracurriculares ou para aprimorar habilidades linguísticas, é necessário fazer isso após o trabalho e o transporte de volta para casa.
Então meu turno começava por volta das 22h30 e terminava 6h30, 7h da manhã.
Muitos estudavam inglês na madrugada.
Esse desejo de aprender era fascinante.
Tanto a experiência no Cazaquistão quanto em Osaka me fizeram ver que estou no campo certo, porque sinto que realmente estou contribuindo com os objetivos educacionais das pessoas.
Depois disso, vivi dois anos na Tailândia ensinando inglês, e outros dois anos na Polônia.
Enfim, senti saudades da Ásia Central e me mudei para o Quirguistão, onde ajudei estudantes quirguizes a estudar nos Estados Unidos.
Finalmente voltei para os Estados Unidos.
Visitando universidades ao redor do mundo, o que poderia dizer sobre Internacionalização e Mobilidade?
Diria que é fascinante ver o crescimento e o interesse [ sobre esses temas ].
Primeiro, devido ao desenvolvimento da tecnologia, crescendo globalmente na ponta dos nossos dedos.
Isso leva a geração mais nova a se interessar e aprender rapidamente, instantaneamente, sobre o que acontece no mundo.
Sinto que a mobilidade é mais rápida, além de mais acessível, porque mesmo que você não consiga viajar fisicamente, você pode aprender sobre outros países assistindo conteúdos online ou fazendo conexões e mantendo melhor os contatos.
Veja como estamos falando agora, eu na Pensilvânia e vocês aí em São Paulo.
Outra questão é que existe mais conhecimento sendo compartilhado.
No tempo em que vivi no Cazaquistão, a não ser que você conhecesse alguém de lá, você não saberia sobre o país a não ser a partir das notícias.
Agora, há outros caminhos.
Por exemplo, a música, que se tornou mais global.
Você vê de fato o impacto da mobilidade global, antes era um presente que apenas poucas pessoas conseguiam.
Agora é mais acessível a todos.
O que considero mais importante é o interesse de pessoas de mente aberta que desejam estar juntas.
Ver paisagens online desperta mais vontade de viajar.
Mas, ainda assim, existe a limitação financeira para viajar fisicamente.
Existe muito medo de viajar a partir de uma pequena cidade para outro lugar do mundo.
Qual sua experiência profissional mais desafiadora e gratificante?
Culturalmente, em todos os lugares em que estive houve uma revelação interessante, mas também um desafio.
Trabalhando com tantas culturas diferentes, há complexidades tão diversas.
Estou constantemente trocando de código.
O que funciona com os brasileiros pode não funcionar com nigerianos ou espanhóis.
Diria que a maioria dos meus desafios profissionais se relaciona com me assegurar de que não esteja sendo ofensiva e reconhecer que a forma como vou me comunicar com um país pode não ser o mesmo método que preciso utilizar com outro.
Gostaria que você comentasse a noção de “global engagement”, que não é muito conhecida aqui no Brasil.
Do meu ponto de vista às vezes a questão é interpretada assim: “quantos estrangeiros entram na universidade e quantos estudantes da instituição saem”.
Mas quando se está trabalhando pelo engajamento global, quais são os outros caminhos?
Adoro as parcerias COIL [ Collaborative Online International Learning ], porque reconhecemos que nem todos terão oportunidades de ir para outro país.
Mas poderão aprender muito virtualmente, por exemplo, sobre as comidas do outro país.
Outra maneira de engajar globalmente é apoiar professores, por exemplo, pesquisando em dois países e áreas diferentes.
Uma outra forma é engajar a comunidade local e despertar a curiosidade sobre compreender outras culturas.
Esse é um bom começo.
Por outro lado, existe muito medo do desconhecido.
Respondendo à sua questão, o conceito de ser engajado globalmente se relaciona com algo feito fisicamente através da sua instituição, mas também por meio de sua curiosidade, de maneira informal.
Que futuro você vislumbra para os projetos COIL na sua instituição e em geral para os próximos cinco anos?
Uma das coisas que vejo é o aumento nas colaborações COIL.
Meu objetivo é criar muitos pontos de contatos entre estudantes e professores por meio de iniciativas COIL.
Também tenho a visão de triangular as colaborações em projetos com mais de dois países.

% matched lemmas: hora
\end{textsample}
